\section{Discussion}
\label{sec:discussion}

First, the two filtering mechanisms are not intrinsically dependent on their order. The output of one stage could be used as the input of the other. We nevertheless chose to apply the main approach first because it is designed as a general filtering mechanism based on the business-rule description and source code. It does not depend on a specific documentation structure and can therefore be applied to different object-oriented projects and programming languages. In contrast, the documentation-based filtering exploits characteristics specific to the Berger-Levrault case, particularly the organization of rules into files and sections and the relationship between their names and source-code class names. For another project, this second stage could therefore rely on different contextual information or could be replaced by another refinement mechanism.

Second, although the evaluation is limited to 100 business rules, these rules cover different parts of the documentation and involve diverse business concepts and vocabularies. They are therefore not concentrated on a single functional area or a homogeneous set of business terms. The variation in the resulting candidate sets also reflects this diversity.

The 96\% recall obtained for RQ1 means that the correct implementation was not retained in four cases. The analysis of these cases shows that the losses originate from R.2, which maps business terms to business entities. Some entities have similar names, and the correct entity may therefore receive a lower similarity score than another, semantically related entity. Since only the highest-ranked entities are retained, the correct entity can be discarded at this stage and cannot be recovered later. This highlights the importance of improving the disambiguation of similar business entities.

The documentation-based filtering introduces two additional false negatives, reducing recall from 96\% to 94\%. These cases are mainly related to differences between documentation terminology and source-code class names. In particular, section names do not always correspond exactly to the names of the classes containing the implementation. Consequently, a correct method may be removed even though its implementation is related to the documented rule.

Overall, the results suggest that the two stages play complementary roles. The main approach provides a general, recall-oriented filtering mechanism, while documentation-based information can further reduce the candidate set when suitable project-specific context is available. This progressive strategy substantially narrows the search space, although its effectiveness depends on the ability to correctly associate business terms and documentation context with source-code entities.